\clearpage\begingroup
\SourceLeaf{129}\HistoricalHead{124}{L}{PHYSIK.\quad NACHLASS.}
\begin{center}[C.]\end{center}
\noindent{\small [S. 309]}
\begin{center}
\textbf{Analyse meines kleinen Ramsdenschen Fernrohrs.}\par\medskip
(Einheit sind Zehntel des Fuss.)
\end{center}
\begin{center}
\renewcommand{\arraystretch}{1.4}\setlength{\tabcolsep}{9pt}
\begin{tabular}{r|c|c|c}
\hline\hline
&[Brennw.]&[Distanzen]&Oeffn[ung]\\
\hline\hline
Obj.&$7{,}866$&\raisebox{-.5\normalbaselineskip}[0pt][0pt]{$8{,}097$}&$0{,}96$\\
{}[Ocul. 1]&$0{,}702$&\raisebox{-.5\normalbaselineskip}[0pt][0pt]{$0{,}857$}&$0{,}35$\\
{}[2]&$0{,}837$&\raisebox{-.5\normalbaselineskip}[0pt][0pt]{$1{,}407$}&$0{,}31$\\
{}[3]&$0{,}842$&\raisebox{-.5\normalbaselineskip}[0pt][0pt]{$1{,}104$}&$0{,}41$\\
{}[4]&$0{,}585$&&$0{,}29$
\end{tabular}\par\medskip
Blendung $10\ \mathrm{mm}=0{,}306$ [Zehntelfuss]
\end{center}
\noindent\rule{\linewidth}{.4pt}
\[
\begin{aligned}
\alpha&=+1{,}0000\\
\alpha'&=+0{,}02937\\
\alpha''&=-0{,}10247\\
\alpha'''&=+0{,}05023\\
\alpha^{\mathrm{IV}}&=+0{,}05661
\end{aligned}\qquad\qquad
\begin{aligned}
\beta&=+0{,}12713\\
\beta'&=-0{,}08530\\
\beta''&=-0{,}03712\\
\beta'''&=+0{,}09678\\
\beta^{\mathrm{IV}}&=0
\end{aligned}
\]
\[
\begin{aligned}
\gamma&=0\\
\gamma'&=+8{,}097\\
\gamma''&=+0{,}931\\
\gamma'''&=-14{,}188\\
\gamma^{\mathrm{IV}}&=+5{,}725
\end{aligned}\qquad\qquad
\begin{aligned}
\delta&=+1\\
\delta'&=+10{,}534\\
\delta''&=-9{,}422\\
\delta'''&=-7{,}666\\
\delta^{\mathrm{IV}}&=+17{,}453.
\end{aligned}
\]
\begin{center}
[Aperturen der Ocularlinsen]\par\medskip
\setlength{\tabcolsep}{1em}\renewcommand{\arraystretch}{1.3}
\begin{tabular}{lc|c}
&wegen Helligkeit&[wegen] Gesichtsfeld\\
I.&$0{,}0282\ [=\omega\alpha']$&$0{,}2491\ [=g\gamma']$\\
II.&$0{,}0984\ [=\omega\alpha'']$&$0{,}0303\ [=g\gamma'']$\\
III.&$0{,}0482\ [=\omega\alpha''']$&$0{,}4366\ [=g\gamma''']$\\
IV.&$0{,}0543\ [=\omega\alpha^{\mathrm{IV}}]$&$0{,}1762\ [=g\gamma^{\mathrm{IV}}].$
\end{tabular}\par\medskip
[Gesichtsfeld $g={}]\ 1^{0}\,45'\,50''$\par\medskip
$[\tfrac12g=3175'']$.
\end{center}

\clearpage\SourceLeaf{130}\HistoricalHead{125}{R}{ANALYSE VON FERNROHREN.}
\begin{center}[Farbiger Rand]\end{center}
\[
\begin{array}{r@{\;}r@{\;}r@{\;}r}
[\alpha'\gamma'f'\ ={}]&+0{,}02937+11{,}534&=&+0{,}3385\\
{}[\alpha''\gamma''f''\ ={}]&-0{,}10247+1{,}112&=&-0{,}1140\\
{}[\alpha'''\gamma'''f'''\ ={}]&+0{,}05023-17{,}088&=&-0{,}8583\\
{}[\alpha^{\mathrm{IV}}\gamma^{\mathrm{IV}}f^{\mathrm{IV}}\ ={}]&+0{,}05661+9{,}787&=&+0{,}5540\\\cline{3-4}
&&=&[-0{,}0798]
\end{array}
\]
\[
\left[-0{,}0798\mathbin{.}\frac{g}{2}=\right]-0{,}0798\mathbin{.}3175''=254''
\]
\[
\left[254''\frac{\partial n}{n-1}=\right]254''\times\frac{0{,}0108}{0{,}515}=5''\text{[*)]}.
\]
\begingroup\renewcommand{\thefootnote}{}\footnotetext{[*) Die Vergrößerung ist auch hier wieder nicht bestimmt; theoretisch soll sie sich gleich dem Werte von $\delta^{\mathrm{IV}}$ ergeben, was im Falle B des \textsc{Fraunhofer}schen Fernrohrs auch wirklich durch den Versuch bestätigt wird.]}\endgroup
\begin{center}\rule{.25\linewidth}{.4pt}\end{center}
\clearpage\endgroup
